% Common-alphabet auxiliaries for SCP10 Theorem 5.5.
% Source: Papers/1001.3807/paper_v3.tex, Definition 5.1 and lines 1421--1513.
% The source-labelled closure theorem is owned by the dependent-alphabet entry.
\subsection{Four cut spaces with common coordinate alphabets}
\label{sec:peps_four_cut_common_boundaries}

The following constructions give the common-coordinate specialization of the
four boundary spaces in \cite[Theorem~5.5]{Schuch2010PEPS}, whose full
statement is Theorem~\ref{thm:peps_four_cut_closure}. Let
$X=\Z/w\Z\times\Z/\ell\Z$, where $w,\ell>0$, and write
$e=(1,0)$ and $n=(0,1)$. Horizontal bonds run from $v$ to $v+e$;
vertical bonds run from $v+n$ to $v$. Let $V$ be a finite virtual alphabet
and $P$ a physical alphabet. All eight bonds remain distinct when
$w=\ell=2$. Write $\mathcal N_A(O^{\mathrm h},O^{\mathrm v})$
for the physical vector obtained from the bond contraction of
Definition~\ref{def:peps_torus_bond_network}, with independent site tensors
$A_v:V^4\times P\to\C$.

\begin{definition}[Correlated cut boundaries and their ranges]%
    \label{def:peps_four_cut_common_boundary}
    \lean{TNLean.PEPS.TorusCutBoundaryConfig,
        TNLean.PEPS.torusCutBoundaryRestriction,
        TNLean.PEPS.torusCutInteriorWeight,
        TNLean.PEPS.torusCutCoeff,
        TNLean.PEPS.torusCutMap,
        TNLean.PEPS.torusCutMap_apply,
        TNLean.PEPS.torusCutSpace,
        TNLean.PEPS.mem_torusCutSpace_iff,
        TNLean.PEPS.fourTorusCutSpace,
        TNLean.PEPS.mem_fourTorusCutSpace_iff,
        TNLean.PEPS.card_torusCutBoundaryConfig_two}
    \leanok
    \uses{def:peps_torus_bond_network}
    The boundary configurations for the cut $(c,r)$ form
    $\mathcal E=(V\times V)^{\Z/\ell\Z}\times
        (V\times V)^{\Z/w\Z}$. Each pair is ordered tail, head.
    For independent bond-end labels $\beta:X\to V^4$, put
    \begin{align}
        R_{c,r}(\beta)
                   & =\bigl(y\mapsto(\beta^1_{(c-1,y)},\beta^2_{(c-1,y)}),
        x\mapsto(\beta^3_{(x,r-1)},\beta^4_{(x,r-1)})\bigr),
        \label{eq:peps_common_boundary_restriction}                        \\
        s_v(\beta) & =(\beta_v^4,\beta_v^1,\beta_{v-n}^3,\beta_{v-e}^2),
        \nonumber                                                          \\
        W_{c,r}(\beta)
                   & =\prod_{v\in X}
        \begin{cases}1,&v_x+1=c,\\\delta_{\beta_v^2,\beta_v^1},&v_x+1\ne c\end{cases}
        \begin{cases}1,&v_y+1=r,\\\delta_{\beta_v^4,\beta_v^3},&v_y+1\ne r.\end{cases}
        \nonumber
    \end{align}
    For an arbitrary joint boundary tensor $M:\mathcal E\to\C$, define
    the linear map $\mathcal C^A_{c,r}$ by
    \begin{align}
        (\mathcal C^A_{c,r}M)(\sigma)
         & =\sum_{\beta:X\to V^4}M(R_{c,r}(\beta))W_{c,r}(\beta)
        \prod_{v\in X}A_v(s_v(\beta);\sigma_v).
        \label{eq:peps_common_cut_contraction}
    \end{align}
    Put $\mathcal S^A_{c,r}=\operatorname{ran}\mathcal C^A_{c,r}$.
    Thus $\psi\in\mathcal S^A_{c,r}$ precisely when some $M$ satisfies
    $(\mathcal C^A_{c,r}M)(\sigma)=\psi(\sigma)$ for every $\sigma$.
    On $X=(\Z/2\Z)^2$, put
    $\mathcal S^A=\bigcap_{c,r}\mathcal S^A_{c,r}$; membership means
    that each cut admits such a boundary tensor, which may depend on the cut.
    Here $|\mathcal E|=|V|^8$.
\end{definition}
\begin{proof}\leanok
    The boundary occurs linearly in (\ref{eq:peps_common_cut_contraction}).
    The range descriptions follow by equality of all coefficients, and the
    intersection description follows by imposing all four range conditions.
    The two rows and two columns each contribute two independent endpoints,
    giving $|V|^{2\cdot2+2\cdot2}=|V|^8$ configurations.
\end{proof}

\begin{definition}[Product boundaries and fixed endpoints]%
    \label{def:peps_four_cut_common_product_boundary}
    \lean{TNLean.PEPS.torusCutBondBoundary,
        TNLean.PEPS.torusCutBondBoundary_restriction,
        TNLean.PEPS.torusCutHorizontalUnit,
        TNLean.PEPS.torusCutVerticalUnit,
        TNLean.PEPS.torusCutBondBoundary_units}
    \leanok
    \uses{def:peps_four_cut_common_boundary}
    For bond matrices $O^{\mathrm h}_v,O^{\mathrm v}_v$, define
    \begin{align}
        M^{O}_{c,r}(\eta)
         & =\prod_y O^{\mathrm h}_{(c-1,y)}
        (\eta^{\mathrm h,2}_y,\eta^{\mathrm h,1}_y)
        \prod_x O^{\mathrm v}_{(x,r-1)}
        (\eta^{\mathrm v,2}_x,\eta^{\mathrm v,1}_x).
        \label{eq:peps_common_product_boundary}
    \end{align}
    Its value at $R_{c,r}(\beta)$ is the product of the corresponding
    $O^{\mathrm h}_v(\beta_v^2,\beta_v^1)$ and
    $O^{\mathrm v}_v(\beta_v^4,\beta_v^3)$ on cut bonds, with factor one
    on every other bond. For fixed $\eta\in\mathcal E$, let
    $E^{\mathrm h}_{\eta,v}=E_{\eta^{\mathrm h,2}_{v_y},
                \eta^{\mathrm h,1}_{v_y}}$ and
    $E^{\mathrm v}_{\eta,v}=E_{\eta^{\mathrm v,2}_{v_x},
                \eta^{\mathrm v,1}_{v_x}}$ be matrix units.
    Then $M^{E_\eta}_{c,r}(\theta)=\delta_{\eta,\theta}$.
\end{definition}
\begin{proof}\leanok
    Substitution of (\ref{eq:peps_common_boundary_restriction}) selects
    the designated seam bonds. Every matrix unit fixes its two endpoints;
    their product is one exactly when all endpoint labels agree with $\eta$.
\end{proof}

\begin{theorem}[Fixed-boundary expansion of a cut range]%
    \label{thm:peps_four_cut_common_boundary_span}
    \lean{TNLean.PEPS.torusCutCoeff_bondBoundary,
        TNLean.PEPS.torusCutMap_eq_sum_single,
        TNLean.PEPS.torusCutSpace_eq_span_single,
        TNLean.PEPS.torusCutMap_single_eq_bondNetwork}
    \leanok
    \uses{def:peps_four_cut_common_product_boundary}
    Write $\widehat O^{\mathrm h,c}_v=O^{\mathrm h}_v$ on
    $v_x+1=c$ and $\Id$ elsewhere, and define
    $\widehat O^{\mathrm v,r}$ analogously. Then
    \begin{align}
        \mathcal C^A_{c,r}M^O_{c,r}
         & =\mathcal N_A(\widehat O^{\mathrm h,c},\widehat O^{\mathrm v,r}),
        \label{eq:peps_common_boundary_network}                              \\
        \mathcal C^A_{c,r}M
         & =\sum_{\eta\in\mathcal E}M(\eta)\mathcal C^A_{c,r}\delta_\eta,
        \nonumber                                                            \\
        \mathcal S^A_{c,r}
         & =\spn_{\C}\{\mathcal C^A_{c,r}\delta_\eta:\eta\in\mathcal E\}.
        \nonumber
    \end{align}
    In particular, $\mathcal C^A_{c,r}\delta_\eta$ is the network with
    $E^{\mathrm h}_{\eta,v}$ and $E^{\mathrm v}_{\eta,v}$ on cut bonds
    and identities elsewhere.
\end{theorem}
\begin{proof}\leanok
    \uses{def:peps_four_cut_common_product_boundary}
    Multiplying the restricted product boundary by $W_{c,r}$ gives exactly
    the bond weights in (\ref{eq:peps_common_boundary_network}). Expand
    $M=\sum_\eta M(\eta)\delta_\eta$ and use linearity. Each fixed-boundary
    vector already belongs to the range, proving the reverse span inclusion.
    The matrix-unit identity gives the final assertion.
\end{proof}

\begin{definition}[Matching bond actions and their averaging tensors]%
    \label{def:peps_four_cut_common_matched_action}
    \lean{TNLean.PEPS.torusMatchedLegMatrix,
        TNLean.PEPS.torusMatchedLegRep,
        TNLean.PEPS.torusMatchedLegRep_apply,
        TNLean.PEPS.torusMatchedLegRep_const,
        TNLean.PEPS.IsUnitaryTorusBondFamily,
        TNLean.PEPS.IsSemiRegularTorusBondFamily,
        TNLean.PEPS.representationAveragingSite,
        TNLean.PEPS.siteMap_representationAveragingSite,
        TNLean.PEPS.representationAveragingSite_torusLegRep,
        TNLean.PEPS.torusMatchedAveragingSites}
    \leanok
    \uses{def:peps_torus_bond_network,def:peps_g_injective,def:peps_semiregular}
    Let $G$ be a group and assign representations $U^{\mathrm h}_v$ and
    $U^{\mathrm v}_v$ on $\C^V$ to the individual bonds. Their action on
    the top, right, down, left coordinates of site $v$ is
    \begin{align}
        \rho_v(k)
         & =U^{\mathrm v}_v(k)\otimes
        U^{\mathrm h}_v(k^{-1})^{\mathsf T}\otimes
        U^{\mathrm v}_{v-n}(k^{-1})^{\mathsf T}\otimes
        U^{\mathrm h}_{v-e}(k).
        \label{eq:peps_common_matched_action}
    \end{align}
    The coefficient action is multiplication by this matrix. The bond
    family is unitary when every $U^d_v(k)$ is unitary, and, for finite
    $G$, is semi-regular when every individual representation $U^d_v$
    is semi-regular. These are separate conditions.
    For any representation $\rho$ on $\C^{V^4}$ of finite $G$, put
    $\Pi_\rho=|G|^{-1}\sum_k\rho(k)$ and
    $A^0_\rho(t,r,b,l;s)=(\Pi_\rho)_{s,(t,r,b,l)}$.
    Its site map is exactly $\Pi_\rho$. Write $A^0_v=A^0_{\rho_v}$
    for the matched family. Constant bond representations recover the
    uniform four-leg action and its averaging tensor.
    The matching convention is that of
    \cite[Definition~5.1]{Schuch2010PEPS}: each link carries its own
    representation, used at both ends.
\end{definition}
\begin{proof}\leanok
    The inverse transpose reverses the order twice, so
    (\ref{eq:peps_common_matched_action}) is a representation.
    The stated coefficient action and the site map of $A^0_\rho$
    follow by reading their matrix entries. Constant bond families
    give the uniform formulas by substitution.
\end{proof}

\begin{definition}[Site-dependent matched closures]%
    \label{def:peps_four_cut_common_closure}
    \lean{TNLean.PEPS.sitewiseTorusGClosure,
        TNLean.PEPS.sitewiseTorusGClosure_const,
        TNLean.PEPS.sitewiseCommutingClosureSpan,
        TNLean.PEPS.torusMatchedHorizontalClosureAt,
        TNLean.PEPS.torusMatchedVerticalClosureAt,
        TNLean.PEPS.matchedTorusGClosure,
        TNLean.PEPS.matchedTorusGClosure_const,
        TNLean.PEPS.matchedCommutingClosureSpan}
    \leanok
    \uses{def:peps_four_cut_common_matched_action,def:peps_torus_g_closure}
    For $g,h\in G$, define
    \begin{align}
        H^c_v(h)          & =\begin{cases}U^{\mathrm h}_v(h),&v_x+1=c,\\\Id,&v_x+1\ne c,\end{cases}
        \nonumber                                                                                   \\
        V^r_v(g)          & =\begin{cases}U^{\mathrm v}_v(g),&v_y+1=r,\\\Id,&v_y+1\ne r,\end{cases}
        \nonumber                                                                                   \\
        \mathcal M_A(g,h) & =\mathcal N_A(H^0(h),V^0(g)),
        \label{eq:peps_common_matched_closure}                                                      \\
        \mathcal Z_A      & =\spn_{\C}\{\mathcal M_A(g,h):gh=hg\}.
        \nonumber
    \end{align}
    When every bond representation is $U$, these are the site-dependent
    uniform closures and their span. If also every site tensor equals $A$,
    the closure is $\mathcal M(A\mid g,h)$ of
    Definition~\ref{def:peps_torus_g_closure}.
\end{definition}

\begin{theorem}[Uniform closures belong to every cut range]%
    \label{thm:peps_four_cut_common_uniform_membership}
    \lean{TNLean.PEPS.torusBondNetwork_closureAt_eq_sitewiseTorusGClosure,
        TNLean.PEPS.torusCutMap_closureBoundary,
        TNLean.PEPS.sitewiseTorusGClosure_mem_torusCutSpace,
        TNLean.PEPS.sitewiseCommutingClosureSpan_le_fourTorusCutSpace}
    \leanok
    \uses{def:peps_four_cut_common_closure,def:peps_four_cut_common_product_boundary}
    Suppose every bond carries $U$ and each site map obeys
    $\mathcal P(A_v)\rho(k)=\mathcal P(A_v)$ for every $k\in G$.
    If $gh=hg$, then for every cut $(c,r)$,
    \begin{align}
        \mathcal N_A(H^c(h),V^r(g))
         & =\mathcal C^A_{c,r}M^{(U(h),U(g))}_{c,r}
        =\mathcal M_A(g,h)\in\mathcal S^A_{c,r}.
        \label{eq:peps_common_uniform_membership}
    \end{align}
    Consequently $\mathcal Z_A\subseteq\mathcal S^A$ on the two-by-two torus.
\end{theorem}
\begin{proof}\leanok
    \uses{thm:peps_torus_commuting_closure_seams,thm:peps_four_cut_common_boundary_span}
    Put $B(t,r,b,l;(v,s))=A_v(t,r,b,l;s)$, with $(v,s)$ as its
    physical label. Its invariance holds separately at each physical label. Apply
    Theorem~\ref{thm:peps_torus_commuting_closure_seams} to $B$ and the
    configuration $v\mapsto(v,\sigma_v)$. The product-boundary identity
    gives the explicit witness in (\ref{eq:peps_common_uniform_membership});
    taking spans proves the inclusion.
\end{proof}

\begin{definition}[Cyclic seam changes of basis]%
    \label{def:peps_four_cut_common_seam_gauge}
    \lean{TNLean.PEPS.torusSeamGauge}
    \leanok
    For $x,c\in\Z/m\Z$, with $m>0$, use representatives
    $\bar x,\bar c\in\{0,\ldots,m-1\}$ and put
    $s_{k,c}(x)=k^{-1}$ if $\bar x<\bar c$, and $s_{k,c}(x)=1$ otherwise.
\end{definition}

\begin{theorem}[Cyclic seam identities]%
    \label{thm:peps_four_cut_common_seam_gauge}
    \lean{TNLean.PEPS.torusSeamGauge_horizontal,
        TNLean.PEPS.torusSeamGauge_vertical,
        TNLean.PEPS.torusSeamGauge_commute}
    \leanok
    \uses{def:peps_four_cut_common_seam_gauge}
    Write $K^c_k(x)=k$ if $x+1=c$, and $K^c_k(x)=1$ otherwise.
    For every $k\in G$,
    \begin{align}
        K^c_k(x) & =s_{k,c}(x+1)K^0_k(x)s_{k,c}(x)^{-1},
        \label{eq:peps_common_seam_horizontal}                     \\
        K^c_k(x) & =s_{k^{-1},c}(x)K^0_k(x)s_{k^{-1},c}(x+1)^{-1}.
        \label{eq:peps_common_seam_vertical}
    \end{align}
    If $kl=lk$, then $s_{k,c}(x)$ commutes with $l$.
\end{theorem}
\begin{proof}\leanok
    Check separately the wrap $x+1=0$, the crossing $x+1=c$, and
    the two intervals between them. The adjacent factors cancel except
    at the seam crossings, proving (\ref{eq:peps_common_seam_horizontal}).
    Invert that formula with $k^{-1}$ in place of $k$ to obtain
    (\ref{eq:peps_common_seam_vertical}). Each $s_{k,c}(x)$ is either
    $1$ or $k^{-1}$, proving commutation.
\end{proof}

\begin{theorem}[Matched commuting closures belong to every cut range]%
    \label{thm:peps_four_cut_common_matched_membership}
    \lean{TNLean.PEPS.torusBondNetwork_matchedClosureAt_eq_matchedTorusGClosure,
        TNLean.PEPS.torusCutMap_matchedClosureBoundary,
        TNLean.PEPS.matchedTorusGClosure_mem_torusCutSpace,
        TNLean.PEPS.matchedCommutingClosureSpan_le_fourTorusCutSpace}
    \leanok
    \uses{def:peps_four_cut_common_closure,def:peps_four_cut_common_product_boundary}
    Suppose $\mathcal P(A_v)\rho_v(k)=\mathcal P(A_v)$ for every $v,k$.
    For $gh=hg$ and every $(c,r)$,
    \begin{align}
        \mathcal N_A(H^c(h),V^r(g))
         & =\mathcal C^A_{c,r}M^{(U^{\mathrm h}(h),U^{\mathrm v}(g))}_{c,r}
        =\mathcal M_A(g,h)\in\mathcal S^A_{c,r}.
        \label{eq:peps_common_matched_membership}
    \end{align}
    On the two-by-two torus, $\mathcal Z_A\subseteq\mathcal S^A$.
    This is the first inclusion in the proof of
    \cite[Theorem~5.5]{Schuch2010PEPS}, with a common coordinate alphabet.
\end{theorem}
\begin{proof}\leanok
    \uses{thm:peps_four_cut_common_seam_gauge,
        thm:peps_four_cut_common_vertex_gauge,thm:peps_four_cut_common_boundary_span}
    Apply the matched vertex-gauge identity first with
    $q_v=s_{h,c}(v_x)$ and then with $q_v=s_{g^{-1},r}(v_y)$.
    Formulas (\ref{eq:peps_common_seam_horizontal}) and
    (\ref{eq:peps_common_seam_vertical}) move the two seams.
    Commutation with the transverse closure label leaves its matrices
    unchanged. Each bond uses its own representation of these group
    identities. The product boundary gives the stated range witness;
    linearity then gives the span inclusion.
\end{proof}
